\section{Upper Limit Setting for the HPS \texorpdfstring{$A'$}{A'} Search}

\subsection{The \CLs{} Method}

\subsubsection{Motivation}

The standard frequentist approach to excluding a signal hypothesis is to
compute a $p$-value under the signal-plus-background hypothesis ($H_1$) and
reject $H_1$ when that $p$-value falls below some threshold $\alpha$ (e.g.\
$\alpha = 0.05$ for 95\% CL).  In regions where the experiment has little
sensitivity — where the expected signal is much smaller than the background
fluctuation — this procedure can spuriously exclude a signal that was never
actually testable.  The \CLs{} method~\cite{Read2002,Junk1999} corrects for
this by normalising the signal $p$-value by the background-only $p$-value,
producing a conservative but well-behaved exclusion criterion.

\subsubsection{Test Statistic}

The HPS limit calculation uses the one-sided profile likelihood ratio test
statistic $\tilde{q}_\mu$, defined as
%
\begin{equation}
  \tilde{q}_\mu =
  \begin{cases}
    -2\ln\dfrac{\mathcal{L}(\mu,\,\hat{\hat{\boldsymbol{\theta}}}_\mu)}
               {\mathcal{L}(\hat{\mu},\,\hat{\boldsymbol{\theta}})}
    & \hat{\mu} \geq 0, \\[10pt]
    -2\ln\dfrac{\mathcal{L}(\mu,\,\hat{\hat{\boldsymbol{\theta}}}_\mu)}
               {\mathcal{L}(0,\,\hat{\hat{\boldsymbol{\theta}}}_0)}
    & \hat{\mu} < 0,
  \end{cases}
\end{equation}
%
where $\mu$ is the signal-strength parameter of interest (POI),
$\boldsymbol{\theta}$ are nuisance parameters encoding systematic
uncertainties, $\hat{\mu}$ and $\hat{\boldsymbol{\theta}}$ are the global
maximum-likelihood estimators, and $\hat{\hat{\boldsymbol{\theta}}}_\mu$ is
the conditional MLE for the nuisance parameters at fixed $\mu$.  The
one-sided form $\tilde{q}_\mu$ ensures that over-fluctuations of the data
relative to the signal hypothesis are not used to exclude the signal.

\subsubsection{\CLs{} Definition}

For a given $\mu$, two $p$-values are computed from the distribution of
$\tilde{q}_\mu$:
%
\begin{itemize}
  \item $p_{s+b}(\mu)$: probability, under $H_1(\mu)$, of observing a test
    statistic at least as signal-like as observed;
  \item $p_b(\mu)$: probability, under $H_0$ (background only), of observing
    a test statistic at least as signal-like.
\end{itemize}
%
The \CLs{} value is then
%
\begin{equation}
  \CLs(\mu) = \frac{p_{s+b}(\mu)}{1 - p_b(\mu)}.
\end{equation}
%
The denominator $1-p_b$ suppresses the exclusion power when the experiment is
insensitive: if the observed data fluctuates toward the background-only
prediction regardless of $\mu$, both numerator and denominator are small and
\CLs{} stays near unity.  A signal hypothesis is excluded at confidence level
CL when $\CLs(\mu) < 1 - \mathrm{CL}$, i.e.\ $\CLs < 0.05$ for 95\% CL.

\subsubsection{Upper Limit}

The 95\% CL upper limit on $\mu$ is the value $\mu_\text{upper}$ where
$\CLs(\mu_\text{upper}) = 0.05$.  In practice this is found by root-finding
over a scan of $\mu$ values.

\subsubsection{Expected Limits and Uncertainty Bands}

To assess the expected sensitivity independently of the observed data, the
calculation is repeated using \emph{Asimov data} — the dataset in which every
observable equals its expected value under the background-only hypothesis.
The median expected limit $\hat{\mu}^\text{exp}$ is the result for Asimov
data.  The $\pm1\sigma$ and $\pm2\sigma$ uncertainty bands are obtained from
the quantiles of $\mu_\text{upper}$ over an ensemble of pseudo-experiments
generated under $H_0$, characterising the range of limits one would expect
from statistical fluctuations alone.

% ============================================================
\subsection{\pyhf{}}
% ============================================================

\pyhf{}~\cite{pyhf,pyhf_joss} is a pure-Python library implementing the
HistFactory statistical model used widely in HEP analyses.  It provides:
%
\begin{itemize}
  \item \textbf{Workspaces} — JSON-serialisable likelihood specifications
    that fully describe a model, its data, and the measurement configuration.
    Workspaces can be published alongside results for reproducibility.

  \item \textbf{Channels and samples} — each \emph{channel} (signal region)
    contains one or more \emph{samples} (e.g.\ signal, background).  Each
    sample carries expected yields and a list of \emph{modifiers} encoding
    systematic uncertainties.

  \item \textbf{Modifiers} — the types used in this framework are:
    \begin{itemize}
      \item \texttt{normfactor}: a freely floating multiplicative scale
        factor; used for the POI $\mu$.
      \item \texttt{normsys}: a log-normal constraint on a sample's
        normalisation, parameterised by the $\pm1\sigma$ scale factors
        $\theta_\text{hi}$ and $\theta_\text{lo}$.
      \item \texttt{shapesys}: a per-bin Poisson-constrained uncertainty,
        parameterised by the absolute yield uncertainty.
    \end{itemize}

  \item \textbf{Inference} — \texttt{pyhf.infer.hypotest} computes \CLs{}
    (and the expected band) for a fixed $\mu$.
    \texttt{pyhf.infer.intervals.upper\_limits.upper\_limit} finds the
    $\mu_\text{upper}$ crossing by root-finding (TOMS~748 algorithm).

  \item \textbf{Backends} — \pyhf{} is backend-agnostic (NumPy, JAX,
    PyTorch), enabling automatic differentiation and GPU acceleration.
\end{itemize}

% ============================================================
\subsection{Implementation in \texttt{hps-limits}}
% ============================================================

\subsubsection{Statistical Model}

For each $A'$ mass hypothesis $m$ and coupling $\varepsilon^2$, the model
contains one \pyhf{} channel per detector hit category (e.g.\ L1L1, L1L2,
L2L2, L2L3, L3L3).  Each channel has two samples:
%
\begin{itemize}
  \item \textbf{signal}: expected yield $N_\text{sig}(m,\varepsilon^2)$
    scaled by the POI $\mu$ via a \texttt{normfactor} modifier;
  \item \textbf{background}: expected yield $N_\text{bkg}(m)$ from the ABCD
    background estimate, with uncertainties described by the modifiers listed
    below.
\end{itemize}
%
The POI $\mu = 1$ corresponds to the nominal signal prediction at the given
$(m, \varepsilon^2)$.  An upper limit $\mu_\text{upper} < 1$ means the
nominal signal is excluded.

\subsubsection{Systematic Uncertainties}

Uncertainties are specified per search region.  Two modifier types are
supported, summarised in Table~\ref{tab:modifiers}.

\begin{table}[h]
  \centering
  \caption{Systematic uncertainty types supported by \texttt{hps-limits}.}
  \label{tab:modifiers}
  \begin{tabular}{lll}
    \hline
    Input type  & \pyhf{} modifier & Interpretation \\
    \hline
    \texttt{normsys}  & \texttt{normsys}  &
      Fractional $1\sigma$ uncertainty $v$;
      $\theta_\text{hi} = 1+v$, $\theta_\text{lo} = 1-v$ \\
    \texttt{shapesys} & \texttt{shapesys} &
      Absolute $1\sigma$ yield uncertainty; Poisson-constrained \\
    \hline
  \end{tabular}
\end{table}

Nuisance parameters with the \emph{same name} across regions are treated as
fully correlated (shared NP in the global fit); distinct names give
uncorrelated NPs.  Background statistical uncertainties are modelled as
uncorrelated \texttt{shapesys} nuisance parameters, one per region.

\subsubsection{Input Data}

\textbf{Background} is loaded from results JSON files produced by the ABCD
background estimator.  Each entry carries $N_\text{bkg}$, its $1\sigma$
uncertainty, and the observed count in the signal window, indexed by (region,
mass).  Luminosity projections can be substituted at load time.

\textbf{Signal} yields are loaded separately from signal or cut-scan JSON
files, providing $N_\text{sig}(\text{region}, m, \varepsilon^2)$ from Monte
Carlo simulation.  Signal is intentionally decoupled from the background
configuration so that the same background model is reused across all mass and
coupling hypotheses.

For $z_0$-binned analyses, the transverse impact parameter $z_0$ is divided
into sub-bins, each becoming a separate \pyhf{} channel, enabling a shape fit
in the displaced-vertex variable.

\subsubsection{Limit-Setting Workflow}

\paragraph{Single mass, signal-strength limit.}
\texttt{compute\_upper\_limit\_from\_input} builds the \pyhf{} model and
calls \texttt{upper\_limit} over a $\mu$ scan of 100 points in
$[0, \mu_\text{max}]$ to find the $\CLs = 0.05$ crossing.

\paragraph{Mass scan.}
\texttt{scan\_mass\_limits\_from\_input} loops over a mapping
$\{m \to \{r \to N_\text{sig}\}\}$, calling
\texttt{compute\_upper\_limit\_from\_input} at each mass point and returning
a list of \texttt{LimitResult} objects.

\paragraph{Direct $\varepsilon^2$ scan.}
For displaced-vertex searches where $N_\text{sig}(\varepsilon^2)$ has a
non-trivial dependence on $\varepsilon^2$, \texttt{scan\_cls\_from\_estimates}
evaluates \texttt{hypotest}($\mu=1$) at each $(m, \varepsilon^2)$ grid point
and interpolates in $\log\varepsilon^2$ to find the $\varepsilon^2$ where
$\CLs = 0.05$.  This approach makes no assumption that
$N_\text{sig} \propto \varepsilon^2$ and directly locates both the lower and
upper boundaries of the exclusion island.

\paragraph{$\mu$ upper-limit grid.}
\texttt{scan\_mu\_upper\_from\_estimates} computes $\mu_\text{upper}$ at each
$(m, \varepsilon^2)$ grid point via root-finding, producing a 2D sensitivity
map.  Points where sensitivity is below a cap ($\mu_\text{upper} \geq
\mu_\text{cap}$) are assigned the cap value with an early-exit optimisation:
once the exclusion island has been fully traversed, remaining high-$\varepsilon^2$
points are filled without calling \pyhf{}.

\subsubsection{Results}

Each mass point yields a \texttt{LimitResult} containing:
%
\begin{itemize}
  \item \texttt{mu\_upper\_obs} / \texttt{mu\_upper\_exp}: observed and
    expected 95\% CL upper limits on $\mu$;
  \item \texttt{mu\_upper\_exp\_band}: expected band at
    $[-2\sigma,\,-1\sigma,\,\text{median},\,+1\sigma,\,+2\sigma]$;
  \item \texttt{eps2\_upper\_obs} / \texttt{eps2\_upper\_exp} /
    \texttt{eps2\_upper\_exp\_band}: the same quantities expressed directly
    in $\varepsilon^2$ (populated by the direct $\varepsilon^2$ scan path);
  \item \texttt{cls\_obs} / \texttt{cls\_exp}: \CLs{} values at $\mu = 1$
    (diagnostic cross-check).
\end{itemize}
%
For the island-of-exclusion visualisation, \texttt{scan\_cls\_grid} returns
the full $\CLs(\varepsilon^2)$ curves per mass, from which both the lower and
upper $\varepsilon^2$ exclusion boundaries are extracted by finding all
$\CLs = 0.05$ sign changes.

% ---- suggested macros (add to your preamble if not already present) ----
% \newcommand{\CLs}{\ensuremath{\mathrm{CL}_s}}
% \newcommand{\pyhf}{\texttt{pyhf}}
%
% ---- suggested bibliography entries -----------------------------------
% @article{Read2002,
%   author  = {Read, A. L.},
%   title   = {Presentation of search results: the {CL}$_s$ technique},
%   journal = {J. Phys. G},
%   volume  = {28}, pages = {2693}, year = {2002},
%   doi     = {10.1088/0954-3899/28/10/313}}
%
% @article{Junk1999,
%   author  = {Junk, T.},
%   title   = {Confidence level computation for combining searches
%              with small statistics},
%   journal = {Nucl. Instrum. Meth. A},
%   volume  = {434}, pages = {435}, year = {1999},
%   doi     = {10.1016/S0168-9002(99)00498-2}}
%
% @software{pyhf,
%   author    = {Lukas Heinrich and Matthew Feickert and Giordon Stark},
%   title     = {pyhf},
%   doi       = {10.5281/zenodo.1169739}}
%
% @article{pyhf_joss,
%   author  = {Heinrich, L. and Feickert, M. and Stark, G. and Cranmer, K.},
%   title   = {pyhf: pure-{P}ython implementation of {HistFactory}},
%   journal = {J. Open Source Softw.},
%   volume  = {6}, number = {58}, pages = {2823}, year = {2021},
%   doi     = {10.21105/joss.02823}}
